\documentclass{article}
\usepackage{booktabs}

\begin{document}

  \section{Punto 1)}
    Repita el análisis para el caso de 16 S-boxes, cada una tomando 
    entradas de 4 bits. ¿Cuál es la complejidad del ataque ahora? Repita 
    el análisis nuevamente con una longitud de bloque de 128 bits y 16 S-boxes
    que cada una toma una entrada de 8 bits.

    \subsection{Análisis para 16 S-boxes con 4 bits}
      Tomemos un par $(x, y)$, donde $y$ es la encriptación de $x$, es decir

        \[y = F_{k_1, k_2}(x) = \pi(S(x \oplus k_1)) \oplus k_2\]

      entonces, el espacio total de claves es $2^{128}$ (inicialmente).
      \newline \newline

      De la ecuación anterior, podemos despejar $k_2 =  \pi(S(x \oplus k_1)) \oplus y$ 
      y notar inmediatamente que la segunda clave está totalmente determinada 
      por $(x, y)$ y $k_1$, entonces el espacio de claves se reduce a $2^{64}$.
      \newline \newline

      Repitiendo el algoritmo de recuperación local, para cada 
      S-box se hace una búsqueda exaustiva en $\{0, 1\}^4$. Cada 
      lista tiene $2^4$ posibilidades de entradas, lo que le costaría
      $16 \cdot 2^4 = 2^8$ operaciones.
      
      A continuación, se presenta el resumen
      \begin{table}[h]
        \centering

        \vspace{0.2cm}
        \begin{tabular}{lcc}
          \toprule
          Etapa & Tamaño & Tiempo \\
          \midrule
          Sin info & $2^{128}$ & --- \\
          Enumerando en $k_1$ & $2^{64}$ & $2^{64}$ \\
          Fragmentado & $2^{64}$ & $2^{8}$ \\
          \bottomrule
        \end{tabular}
      \end{table}

    \subsection{Análisis para 16 S-boxes con 8 bits}
      Tomemos un par $(x, y)$, donde $y$ es la encriptación de $x$, es decir

        \[y = F_{k_1, k_2}(x) = \pi(S(x \oplus k_1)) \oplus k_2\]

      entonces, el espacio total de claves es $2^{256}$ (inicialmente).
      \newline \newline

      De la ecuación anterior, podemos despejar $k_2 =  \pi(S(x \oplus k_1)) \oplus y$ 
      y notar inmediatamente que la segunda clave está totalmente determinada 
      por $(x, y)$ y $k_1$, entonces el espacio de claves se reduce a $2^{128}$.
      \newline \newline

      Repitiendo el algoritmo de recuperación local, para cada 
      S-box se hace una búsqueda exaustiva en $\{0, 1\}^8$. Cada 
      lista tiene $2^8$ posibilidades de entradas, lo que le costaría
      $16 \cdot 2^8 = 2^{12}$ operaciones.

      A continuación, se presenta el resumen
      \begin{table}[h]
        \centering

        \vspace{0.2cm}
        \begin{tabular}{lcc}
          \toprule
          Etapa & Tamaño & Tiempo \\
          \midrule
          Sin info & $2^{256}$ & --- \\
          Enumerando en $k_1$ & $2^{128}$ & $2^{128}$ \\
          Fragmentado & $2^{128}$ & $2^{12}$ \\
          \bottomrule
        \end{tabular}
      \end{table}

  \section{Punto 2)}
    Considere un SPN de dos rondas, el cifrado es de la forma
      \[F_{k_1, k_2, k_3}(x) = \pi(S(\pi(S(x \oplus k_1)) \oplus k_2)) \oplus k_3\]

    Haga el análisis con SPN de dos rondas con bloques de 64 bits y 
    S-boxes de 8 bits y construya un ataque que recupere las claves 
    en un tiempo menor a $2^{192}$.

    Para empezar, vemos que inicialmente, dado un par $(x, y)$, con 
      \[y = \pi(S(\pi(S(x \oplus k_1)) \oplus k_2)) \oplus k_3\]

    el tamaño de las claves es $2^{192}$ posibles claves.

    De la ecuación del enunciado, vemos que 
      \[k_3 = \pi(S(\pi(S(x \oplus k_1)) \oplus k_2)) \oplus y\]


    El análisis se vuelve bastante complejo al tratar de pensar 
    cómo actuaría si se itera en forma de bloque en la primera ronda de S-boxes,
    entonces construimos el siguiente algoritmo:
    \newline \newline

    Para cada posible clave $b \in \{0, 1\}^{64}$, hacemos 
    $w = \pi(S(x \oplus k_1))$ y vemos que 
      \[y = \pi(S(w \oplus k_2)) \oplus k_3\]

    Y una vez tenemos esto, aplicamos el ataque de 8 S-boxes con 
    $8$ bits cada una, para el cual, en el mejor de los casos, resolviamos en 
    tiempo, $8 \cdot 2^8 = 2^{11}$, entonces, el tiempo de resolución de esta forma es 
    $2^{64} \cdot 2^{11} = 2^{75}$ con un tamaño de claves $2^{128}$, pues 
    son $2^{64}$ posibles iniciales y son $2^{64}$ para cada búsqueda.

    \begin{table}[h]
      \centering

      \vspace{0.2cm}
      \begin{tabular}{lcc}
        \toprule
        Etapa & Tamaño & Tiempo \\
        \midrule
        Sin info & $2^{192}$ & --- \\
        Enumerando en $k_1$ y $k_2$ & $2^{128}$ & $2^{128}$ \\
        Fragmentado & $2^{128}$ & $2^{75}$ \\
        \bottomrule
      \end{tabular}
    \end{table}

\end{document}
